\documentclass[12pt,a4paper]{amsart}

\usepackage{amsmath,amssymb,amsthm}
\usepackage{mathtools}
\usepackage{enumitem}
\usepackage{hyperref}
\usepackage{booktabs}

%%%--- Theorem environments
\newtheorem{theorem}{Theorem}[section]
\newtheorem{lemma}[theorem]{Lemma}
\newtheorem{proposition}[theorem]{Proposition}
\newtheorem{corollary}[theorem]{Corollary}
\theoremstyle{definition}
\newtheorem{definition}[theorem]{Definition}
\newtheorem{problem}[theorem]{Open Problem}
\theoremstyle{remark}
\newtheorem{remark}[theorem]{Remark}
\newtheorem*{remark*}{Remark}

%%%--- Macros
\newcommand{\norm}[1]{\|#1\|}
\newcommand{\ip}[2]{\langle #1,\, #2\rangle}
\newcommand{\Hc}{\mathcal{H}_c}
\newcommand{\Hspec}{\mathcal{H}^{\mathrm{spec}}}
\newcommand{\HSOT}{\mathcal{H}^{\mathrm{SOT}}}
\newcommand{\Dom}{\mathrm{Dom}}

\subjclass[2020]{47B25, 47A10, 42C05, 34E20, 11M06}
\keywords{Prolate spheroidal wave functions, limit operator,
operator closure, essential self-adjointness, open problems}

\begin{document}

\title[PSWF Limit Operator: A Conditional Framework]{%
  The PSWF Concentration Limit Operator:\\%
  A Conditional Framework for Domain,
  Closability, and Self-Adjointness}

\author{[Author]}
\address{[Address]}
\email{[Email]}
\maketitle

\begin{abstract}
We develop a conditional functional-analytic framework
for the spectral study of the PSWF concentration
operator family $\{\Hc\}$ and its two natural
limit objects:
the SOT-limit $\HSOT$ (a bounded self-adjoint operator)
and the spectral operator $\Hspec$ (defined by a
spectral series over subsequential limit vectors).

This paper does \emph{not} prove that $\Hspec$
is essentially self-adjoint, that $\Hspec$
and $\HSOT$ are unitarily equivalent, or any
statement about the Riemann zeta function.
Instead, it identifies precisely which
hypotheses would be needed for such conclusions
and what is unconditionally established.

\textbf{Unconditional results:}
precompactness of $\{\Phi_n^{(c)}\}$ in $L^2$;
closability and symmetry of $\Hspec$;
boundedness and self-adjointness of $\HSOT$.

\textbf{Conditional results} (explicit hypotheses
listed in Sec.~\ref{sec:hypotheses}):
norm convergence $\|\Phi_n^{(c)}-\Phi_n^{(\infty)}\|\to 0$;
the eigenrelation $\HSOT\Phi_n^{(\infty)}
=\lambda_n^{(\infty)}\Phi_n^{(\infty)}$;
essential self-adjointness of $\Hspec$.

\textbf{Central open problems}
(Sec.~\ref{sec:problems}):
norm-resolvent convergence;
discrete spectrum of $\HSOT$;
Riesz-basis property of $\{\Phi_n^{(\infty)}\}$;
completeness.
\end{abstract}

\tableofcontents

%=============================================================
\section{Setup and Inherited Results}
\label{sec:setup}
%=============================================================

We work in $H:=L^2(\mathbb{R})$.
The PSWF concentration operators $\Hc$ are bounded,
self-adjoint, $\|\Hc\|\le 1$, with simple positive
eigenvalues $\lambda_0^{(c)}>\lambda_1^{(c)}>\cdots>0$
and ONS of eigenfunctions $\{\Phi_n^{(c)}\}_n$.

Inherited:
\begin{enumerate}[(P1)]
\item\label{P1}
  $\|(\Hc-\HSOT)f\|\le C_\kappa c^{-3/4}\|f\|$
  for each fixed unit vector $f$
  (Paper~VII, Thm.~C).
\item\label{P2}
  $|\lambda_n^{(c)}-\lambda_n^{(\infty)}|\le C_\kappa c^{-1/4}$,
  $\lambda_n^{(\infty)}:=\lim_c\lambda_n^{(c)}>0$
  (Paper~VIII, Cor.~4.3).
\item\label{P3}
  $\Hc\xrightarrow{\mathrm{SOT}}\HSOT$;
  $\HSOT$ is bounded self-adjoint, $\|\HSOT\|\le 1$
  (Paper~X, Lem.~2.2; see also
  Rem.~\ref{rem:HSOT-bounded}).
\end{enumerate}

\begin{remark}[$\HSOT$ is a bounded self-adjoint operator]
\label{rem:HSOT-bounded}
The SOT-limit of a uniformly bounded net of
bounded self-adjoint operators is bounded self-adjoint
(\cite{ReedSimon1975}, Problem II.3).
With $\|\Hc\|\le 1$ for all~$c$, $\HSOT$ satisfies
$\Dom(\HSOT)=H$ and $\|\HSOT\|\le 1$.
Every bounded everywhere-defined operator is closable;
no separate closability argument is needed.
\end{remark}

%=============================================================
\section{Hypotheses}
\label{sec:hypotheses}
%=============================================================

We list all assumptions on which conditional results
depend. None is proved in this paper.

\begin{hypothesis}[Discrete spectrum of $\HSOT$]
\label{hyp:HSOT-discrete}
$\HSOT$ has purely discrete spectrum with simple
eigenvalues $\lambda_0^{(\infty)}>\lambda_1^{(\infty)}
>\cdots>0$ and corresponding unit eigenvectors
$\{\Phi_n^{(\infty)}\}$.
\end{hypothesis}

\begin{remark}[Spectral decomposition instability under SOT convergence]
\label{rem:HSOT-discrete-hard}
SOT convergence of $\Hc$ to $\HSOT$ does not in general
imply stability of the point spectrum.
Even for uniformly bounded self-adjoint families,
eigenvalues can dissolve into continuous spectrum
in the limit (\cite{ReedSimon1975}, Ex.~VIII.3.4).
We refer to this phenomenon as
\emph{spectral decomposition instability under SOT convergence}.
That the Slepian ordering $\lambda_n^{(\infty)}
>\lambda_{n+1}^{(\infty)}>0$ persists follows from
\hyperref[P2]{(P2)} and the finite-$c$ gap,
but this does not address completeness of the eigensystem.
The problem of restoring stability is central in this work
and is formalized as Hyp.~\ref{hyp:HSOT-discrete}.
Hyp.~\ref{hyp:HSOT-discrete} is the central
``black box'' of the conditional theory.
\end{remark}

\begin{hypothesis}[Norm-resolvent convergence]
\label{hyp:norm-resolvent}
For some (equivalently every)
$z\in\mathbb{C}\setminus\mathbb{R}$:
$\|(\Hc-z)^{-1}-(\HSOT-z)^{-1}\|_{\mathcal{B}(H)}
\to 0$.
\end{hypothesis}

\begin{remark}[SOT $\not\Rightarrow$ norm-resolvent]
\label{rem:SOT-not-enough}
SOT convergence \hyperref[P3]{(P3)} does not imply
norm-resolvent convergence;
see \cite{ReedSimon1975}, Ex.~VIII.7.14.
\end{remark}

\begin{hypothesis}[Completeness]
\label{hyp:decay}
$\{\Phi_n^{(\infty)}\}$ is complete in $L^2(\mathbb{R})$.
\end{hypothesis}

\begin{hypothesis}[Uniform alignment]
\label{hyp:uniform-alignment}
For each $\kappa>0$:
$\sup_{n\le\kappa c}\|\Phi_n^{(c)}-\Phi_n^{(\infty)}\|
\le C_\kappa c^{-1/4}$.
\end{hypothesis}

\begin{remark}[Dependency diagram]
\label{rem:dep}
The conditional results chain as follows:
\[
  \underbrace{\text{Hyp.~\ref{hyp:norm-resolvent}}
    +\text{Hyp.~\ref{hyp:HSOT-discrete}}}
  _{\text{open}}
  \Longrightarrow
  \text{(C1: full conv.)}
  \xRightarrow{+\,\text{Hyp.~\ref{hyp:decay}}}
  \text{(C2: ESSA)}.
\]
Hyp.~\ref{hyp:uniform-alignment} conjecturally implies
Hyp.~\ref{hyp:norm-resolvent} under an additional
Riesz-basis stability condition
(Rem.~\ref{rem:opnorm-bridge}), but this is not proved.
\end{remark}

%=============================================================
\section{Unconditional Results}
\label{sec:unconditional}
%=============================================================

\subsection{Precompactness (no operator)}

\begin{lemma}[PSWF precompactness]
\label{lem:pswf-precompact}
For fixed $n\ge 0$, $\{\Phi_n^{(c)}\}_{c\ge 1}$
is precompact in $L^2$.
Every subsequence has a further subsequence converging
strongly to a unit vector $u\in L^2$.
\end{lemma}

\begin{proof}
Fr\'echet--Kolmogorov (\cite{ReedSimon1975}, Thm.~IV.26):
(i)~$\|\Phi_n^{(c)}\|=1$;
(ii)~uniform exponential decay (\cite{Osipov2013},
Thm.~3.2);
(iii)~uniform $H^1$-bound (\cite{Osipov2013},
Prop.~4.1).
\end{proof}

\begin{remark}
Precompactness guarantees existence of accumulation
points with unit norm but does \emph{not} imply their
uniqueness.
\end{remark}

\subsection{Rank-1 structure at finite $c$}

\begin{lemma}
\label{lem:rank1}
For each finite $c$, the Slepian gap
(\cite{Slepian1961}) ensures:
$P_n^{(c)}f=\ip{f}{\Phi_n^{(c)}}\Phi_n^{(c)}$,
$\|P_n^{(c)}\|=1$.
No uniformity in $c$ is claimed.
\end{lemma}

\subsection{The spectral operator $\Hspec$}

\begin{definition}[$\Hspec$]
\label{def:domain}
Fix any ONS $\{\Phi_n^{(\infty)}\}\subset L^2$
(e.g.\ chosen as subsequential limits;
well-definedness requires the ONS property only).
\begin{align*}
  \Dom(\Hspec)&:=\Bigl\{f\in L^2:\,
  {\textstyle\sum_n}(\lambda_n^{(\infty)})^2
  |\ip{f}{\Phi_n^{(\infty)}}|^2<\infty\Bigr\},\\
  \Hspec f&:=\sum_n\lambda_n^{(\infty)}
  \ip{f}{\Phi_n^{(\infty)}}\Phi_n^{(\infty)}.
\end{align*}
\end{definition}

\begin{remark}[Domain density]
\label{rem:domain-density}
$\mathrm{span}\{\Phi_n^{(\infty)}\}\subset\Dom(\Hspec)$
(direct verification).
Under Hyp.~\ref{hyp:decay} this span is dense,
so $\Dom(\Hspec)$ is dense.
Without Hyp.~\ref{hyp:decay}, density of
$\Dom(\Hspec)$ is \emph{not guaranteed};
however, closability and symmetry do not require
density of the domain in the following proof.
\end{remark}

\begin{theorem}[Closability and symmetry of $\Hspec$]
\label{thm:closable-spec}
$\Hspec$ is closable and symmetric.
\end{theorem}

\begin{proof}
\textit{Symmetry:}
$\ip{\Hspec f}{g}
=\sum_n\lambda_n^{(\infty)}
\ip{f}{\Phi_n^{(\infty)}}\overline{\ip{g}{\Phi_n^{(\infty)}}}
=\ip{f}{\Hspec g}$
for $f,g\in\Dom(\Hspec)$;
the sum converges absolutely by
Cauchy--Schwarz and the $\ell^2$-membership condition.

\textit{Closability:}
Let $U\colon L^2(\mathbb{R})\to\ell^2(\mathbb{N})$,
$Uf:=(\ip{f}{\Phi_n^{(\infty)}})_n$,
be the unitary map (isometric embedding) associated
with the fixed orthonormal system $\{\Phi_n^{(\infty)}\}$.
Let $f_k\to 0$ in $L^2$ and $\Hspec f_k\to g$.
Since $U$ is an isometry,
$f_k\to 0$ in $L^2$ implies
$Uf_k\to 0$ in $\ell^2$,
i.e.\ for each $n$:
$a_n^{(k)}:=\ip{f_k}{\Phi_n^{(\infty)}}\to 0$.
For any $\phi\in L^2$:
\[
  \ip{\Hspec f_k}{\phi}
  =\sum_n\lambda_n^{(\infty)}
  a_n^{(k)}\ip{\Phi_n^{(\infty)}}{\phi}.
\]
The interchange of limit and sum is justified by
dominated convergence in $\ell^2$:
the dominating sequence
$\bigl(\lambda_n^{(\infty)}|\ip{\Phi_n^{(\infty)}}{\phi}|\bigr)_n$
lies in $\ell^2$
(since $\{\lambda_n^{(\infty)}\}$ is bounded and
$\phi\in L^2$ gives $(\ip{\Phi_n^{(\infty)}}{\phi})_n\in\ell^2$
by Bessel's inequality),
and $a_n^{(k)}\to 0$ for each $n$.
Hence $\ip{\Hspec f_k}{\phi}\to 0$ for all $\phi$,
so $g=0$.
\end{proof}

%=============================================================
\section{Conditional: Full Eigenfunction Convergence}
\label{sec:phi-conv}
%=============================================================

\begin{theorem}[Full norm convergence]
\label{thm:phi-conv}
\emph{Conditional on
Hyp.~\ref{hyp:norm-resolvent} and
Hyp.~\ref{hyp:HSOT-discrete}.}
$\|\Phi_n^{(c)}-\Phi_n^{(\infty)}\|\to 0$.
\end{theorem}

\begin{proof}
\textbf{Step~1.}
Hyp.~\ref{hyp:HSOT-discrete}: $\lambda_n^{(\infty)}$
is simple isolated in $\sigma(\HSOT)$;
eigenprojection $P_n^{(\infty)}$ has rank~1.

\textbf{Step~2.}
Kato (\cite{Kato1966}, Thm.~VIII.1.8) under
Hyp.~\ref{hyp:norm-resolvent}:
$\|P_n^{(c)}-P_n^{(\infty)}\|_{\mathcal{B}(H)}\to 0$.

\textbf{Step~3.}
Both $P_n^{(c)},P_n^{(\infty)}$ have rank~1.
From projection-norm convergence:
$|\ip{\Phi_n^{(c)}}{\Phi_n^{(\infty)}}|^2
=\mathrm{tr}(P_n^{(c)}P_n^{(\infty)})\to
\mathrm{tr}(P_n^{(\infty)})=1$,
so $\|\Phi_n^{(c)}-\Phi_n^{(\infty)}\|^2
=2-2\,\mathrm{Re}\ip{\Phi_n^{(c)}}{\Phi_n^{(\infty)}}\to 0$
(positive phase convention).
\end{proof}

%=============================================================
\section{Conditional: Eigenrelation for $\HSOT$}
\label{sec:eigenrelation}
%=============================================================

\begin{lemma}[$\HSOT$-eigenrelation]
\label{lem:weak-eigenrelation}
\emph{Conditional on Thm.~\ref{thm:phi-conv}.}
$\HSOT\Phi_n^{(\infty)}=\lambda_n^{(\infty)}\Phi_n^{(\infty)}$.
\end{lemma}

\begin{proof}
Decompose:
\begin{equation}\label{eq:decompose}
  \HSOT\Phi_n^{(c)}
  = \underbrace{\Hc\Phi_n^{(c)}}_{=\,\lambda_n^{(c)}\Phi_n^{(c)}}
  + \underbrace{(\HSOT-\Hc)\Phi_n^{(c)}}_{=:\,r_c}.
\end{equation}
For the remainder: by \hyperref[P1]{(P1)} with $f=\Phi_n^{(c)}$
and $\|\Phi_n^{(c)}\|=1$:
\begin{equation}\label{eq:remainder}
  \|r_c\| = \|(\HSOT-\Hc)\Phi_n^{(c)}\|
  \le C_\kappa c^{-3/4} \to 0.
\end{equation}
\hyperref[P1]{(P1)} applies here because it holds
\emph{for each fixed unit vector};
applied to the specific unit vector $\Phi_n^{(c)}$,
it gives a bound uniform in the fixed-$n$ sense
(not requiring uniformity over $n$ or $c$ simultaneously).

By Thm.~\ref{thm:phi-conv},
$\|\Phi_n^{(c)}-\Phi_n^{(\infty)}\|\to 0$, so:
\begin{align}
  \HSOT\Phi_n^{(c)}&\to\HSOT\Phi_n^{(\infty)}
  \label{eq:HSOT-norm}
  && (\text{boundedness of }\HSOT,
     \text{ Rem.~\ref{rem:HSOT-bounded}}),\\
  \lambda_n^{(c)}\Phi_n^{(c)}&\to
  \lambda_n^{(\infty)}\Phi_n^{(\infty)}
  \label{eq:eigenvec-limit}
  && (\hyperref[P2]{(P2)}\text{ and Thm.~\ref{thm:phi-conv}}).
\end{align}
Combining \eqref{eq:decompose}--\eqref{eq:eigenvec-limit}:
$\HSOT\Phi_n^{(\infty)}
=\lambda_n^{(\infty)}\Phi_n^{(\infty)}$.

\medskip
\noindent\textit{Note on the mixed error term.}
The key step \eqref{eq:remainder} bounds
$\|(\HSOT-\Hc)\Phi_n^{(c)}\|$ by $C_\kappa c^{-3/4}$
using $\|\Phi_n^{(c)}\|=1$.
This is the only place where $\Phi_n^{(c)}$ (not
$\Phi_n^{(\infty)}$) appears inside a limit;
it is fully justified by \hyperref[P1]{(P1)}
because \hyperref[P1]{(P1)} holds for all unit vectors,
and $\Phi_n^{(c)}$ is a unit vector for each~$c$.
No interchange of operator and vector limits
beyond \eqref{eq:HSOT-norm} is used,
and \eqref{eq:HSOT-norm} is justified by
boundedness of $\HSOT$ alone.
\end{proof}

\begin{remark}[The one genuine approximation]
\label{rem:genuine-approx}
The bound \eqref{eq:remainder} uses \hyperref[P1]{(P1)},
which is a pointwise-in-$f$ estimate.
Applied to $f=\Phi_n^{(c)}$ (which varies with~$c$),
it gives a $c$-dependent bound.
This is valid: for each fixed $c$,
$\|r_c\|\le C_\kappa c^{-3/4}\to 0$.
There is no hidden uniformity over $n$ assumed.
\end{remark}

%=============================================================
\section{Conditional: Essential Self-Adjointness}
\label{sec:esa}
%=============================================================

\begin{theorem}[Completeness implies ESSA]
\label{thm:esa}
\emph{Conditional on
Thm.~\ref{thm:phi-conv} and Hyp.~\ref{hyp:decay}.}
If $\{\Phi_n^{(\infty)}\}$ is complete, then $\Hspec$
is essentially self-adjoint.
\end{theorem}

\begin{proof}
Let $f\in\ker(\Hspec^*\pm i)$.
For each $n$, since $\Phi_n^{(\infty)}\in\Dom(\Hspec)$:
$(\lambda_n^{(\infty)}\pm i)\ip{f}{\Phi_n^{(\infty)}}=0$,
hence $\ip{f}{\Phi_n^{(\infty)}}=0$.
Hyp.~\ref{hyp:decay} gives $f=0$; so $n_\pm=0$.
\end{proof}

\begin{remark}[One direction only; converse is open]
\label{rem:essa-one-direction}
Thm.~\ref{thm:esa} gives only completeness
$\Rightarrow$ ESSA.
The converse requires an explicit description
of $\Dom(\Hspec^*)$, which is not available;
see Prob.~\ref{prob:essa-converse}.
\end{remark}

\begin{corollary}[Under completeness]
\label{cor:esa}
\emph{Conditional as in Thm.~\ref{thm:esa}.}
$\overline{\Hspec}=\Hspec^*$.
\end{corollary}

%=============================================================
\section{Conditional: Operator Norm Rate}
\label{sec:opnorm}
%=============================================================

\begin{theorem}[Partial spectral truncation estimate]
\label{thm:opnorm}
\emph{Conditional on Hyp.~\ref{hyp:uniform-alignment}.}
$\|\Hc-\Hspec\|\le C_\kappa c^{-1/4}$.

\noindent
This estimate controls the spectral truncation error
for $n\le\kappa c$ only and does \emph{not} imply
norm convergence of $\Hc$ to $\HSOT$
in $\mathcal{B}(L^2)$.
\end{theorem}

\begin{remark}[Missing bridge to norm-resolvent]
\label{rem:opnorm-bridge}
Hyp.~\ref{hyp:uniform-alignment} controls
$\|\Phi_n^{(c)}-\Phi_n^{(\infty)}\|$
for $n\le\kappa c$ only.
Operator-norm convergence $\|\Hc-\HSOT\|\to 0$
(which would imply Hyp.~\ref{hyp:norm-resolvent}
via \cite{Kato1966}, Thm.~VIII.1.5)
would further require:
\begin{enumerate}[\rm(i)]
\item negligibility of the tail $n>\kappa c$
  uniformly in $c$; and
\item uniform frame bounds for $\{\Phi_n^{(\infty)}\}$
  (e.g.\ Riesz-basis property;
  see Prob.~\ref{prob:riesz}).
\end{enumerate}
Neither (i) nor (ii) is established here.
Thm.~\ref{thm:opnorm} is therefore a partial
spectral truncation estimate for $\|\Hc-\Hspec\|$,
not for $\|\Hc-\HSOT\|$,
and does not close the bridge to
Hyp.~\ref{hyp:norm-resolvent}.
\end{remark}

%=============================================================
\section{Meta-Remark: Zeta Connection}
\label{sec:rh-remark}
%=============================================================

\begin{remark}[What a genuine zeta theorem would require]
\label{rem:rh-meta}
A spectral theorem connecting $\Hspec$ to
nontrivial zeros of $\zeta$ would need, at minimum:
\begin{enumerate}[(a)]
\item Proof of Hyp.~\ref{hyp:decay} (completeness).
\item Proof of Hyp.~\ref{hyp:norm-resolvent}.
\item A spectral \emph{isomorphism} between
  $\Hspec$ and a trace-class operator whose
  eigenvalues biject with zeta zeros (not just
  an approximate-spectrum inclusion).
\item A Weil-type explicit formula;
  CCM~\cite{CCM2025} does not supply this in a
  form directly applicable here.
\end{enumerate}
None of (a)--(d) is provided in this paper.
The present work identifies the framework and the gaps;
it does not prove RH or any equivalent statement.
\end{remark}

%=============================================================
\section{Open Problems}
\label{sec:problems}
%=============================================================

\begin{problem}[Norm-resolvent convergence]
\label{prob:norm-resolvent}
Prove Hyp.~\ref{hyp:norm-resolvent}.
\end{problem}

\begin{problem}[Discrete spectrum of $\HSOT$]
\label{prob:HSOT-discrete}
Prove Hyp.~\ref{hyp:HSOT-discrete}.
This is the central open problem;
see Rem.~\ref{rem:HSOT-discrete-hard}.
\end{problem}

\begin{problem}[Completeness]
\label{prob:completeness}
Prove Hyp.~\ref{hyp:decay}.
\end{problem}

\begin{problem}[Uniform alignment]
\label{prob:uniform-alignment}
Prove Hyp.~\ref{hyp:uniform-alignment}.
\end{problem}

\begin{problem}[Riesz-basis property]
\label{prob:riesz}
Do the vectors $\{\Phi_n^{(\infty)}\}$ form a Riesz
basis in $L^2$?
This would close gap (i)--(ii) in
Rem.~\ref{rem:opnorm-bridge}.
\end{problem}

\begin{problem}[Converse ESSA]
\label{prob:essa-converse}
Does ESSA imply completeness?
Requires explicit $\Dom(\Hspec^*)$.
\end{problem}

\begin{problem}[Operator identification]
\label{prob:identification}
Prove $\overline{\Hspec}=\HSOT$.
\end{problem}

\begin{thebibliography}{99}
\bibitem{Slepian1961}
  D.~Slepian, H.~O.~Pollak,
  \textit{Prolate spheroidal wave functions~I},
  Bell Syst.\ Tech.\ J.\ \textbf{40} (1961), 43--63.
\bibitem{Osipov2013}
  A.~Osipov, V.~Rokhlin, H.~Xiao,
  \textit{Prolate Spheroidal Wave Functions of Order Zero},
  Springer, 2013.
\bibitem{Kato1966}
  T.~Kato,
  \textit{Perturbation Theory for Linear Operators},
  Springer, 1966.
\bibitem{ReedSimon1975}
  M.~Reed, B.~Simon,
  \textit{Methods of Modern Mathematical Physics, Vol.~II},
  Academic Press, 1975.
\bibitem{CCM2025}
  A.~Connes, C.~Consani, H.~Moscovici,
  \textit{Zeta Spectral Triples}, arXiv:2511.22755, 2025.
\end{thebibliography}

\end{document}
